% Source: Wald GR, physical PDF p. 242 (printed p. 229), Figure 9.3.
% Coverage: Timelike variation after a conjugate point.
% Figures/tables/footnotes: Figure 9.3.
% Status: source-reviewed and compile-clean.
% Uncertainties: none.

\begingroup
\renewcommand{\thefigure}{9.3}
\begin{figure}[tb]
  \centering
  \begin{tikzpicture}[x=1cm,y=1cm,>=Stealth,thick]
    \coordinate (p) at (0,0);
    \coordinate (r) at (.28,2.35);
    \coordinate (q) at (.72,4.2);
    % The varied curve follows a different route only up to the conjugate
    % point r; above r it coincides with gamma (shown dashed).
    \draw (p) .. controls (.02,1.30) and (.36,3.15) .. (q);
    \draw (p) .. controls (-.62,1.48) and (-.36,2.18) .. (r);
    \draw[dashed] (r) -- (q);
    \fill (p) circle (2pt) node[below] {\(p\)};
    \fill (r) circle (2pt) node[right] {\(r\)};
    \fill (q) circle (2pt) node[above right] {\(q\)};
    \draw[->] (1.20,2.12) to[out=155,in=0] (.40,2.08);
    \node[right] at (1.12,1.72) {\(\gamma\)};
    \draw[->] (-1.22,1.48) to[out=45,in=180] (-.31,2.04);
    \node[left] at (-1.1,1.24) {\(\gamma'\)};
  \end{tikzpicture}
  \caption{A spacetime diagram illustrating the fact that if a timelike geodesic
  \(\gamma\) connecting points \(p\) and \(q\) has a point \(r\) conjugate to \(p\) lying
  between \(p\) and \(q\), then a nearby timelike curve \(\gamma'\) connecting \(p\)
  and \(q\) can be constructed which has greater elapsed proper time than \(\gamma\)
  (see theorem~9.3.3).}
  \label{fig:9.3}
\end{figure}
\endgroup
